\section{A Complete Bulletproofs-Style Polynomial Opening Example}

This appendix gives a concrete exponent-level example of the Bulletproofs / IPA-style folding idea \cite{Bulletproofs2018}. The example is not cryptographically secure because the field is tiny; it is meant to make the algebra visible.

\subsection{The polynomial opening claim}

Work over
\[
    \F_{97}.
\]
Let
\[
    f(X)=2+4X+X^2+3X^3.
\]
The coefficient vector is
\[
    \mathbf f=(2,4,1,3).
\]
Let the evaluation point be
\[
    u=5.
\]
Then
\[
\begin{aligned}
    v=f(5)
    &=2+4\cdot 5+5^2+3\cdot 5^3\\
    &=2+20+25+375\\
    &=422\\
    &\equiv 34\pmod{97}.
\end{aligned}
\]
Equivalently,
\[
    v=\iprod{(2,4,1,3)}{(1,5,25,28)}=34,
\]
since $5^3=125\equiv 28\pmod{97}$.

\subsection{Pedersen vector commitment}

Let
\[
    \pp=(g_0,g_1,g_2,g_3)
\]
be four independent group bases in a cyclic group of order $97$. The commitment is
\[
    \mathsf{com}_f=g_0^2g_1^4g_2^1g_3^3.
\]
The prover wants to prove that this committed vector evaluates to $34$ at $u=5$.

\subsection{First folding round}

Split the polynomial into low and high halves:
\[
    \mathbf f_L=(2,4),
    \qquad
    \mathbf f_R=(1,3).
\]
Compute
\[
    v_L=2+4\cdot 5=22,
\]
\[
    v_R=1+3\cdot 5=16.
\]
Then
\[
    v_L+u^2v_R=22+25\cdot 16=422\equiv 34=v.
\]

The prover sends
\[
    L_1=g_2^2g_3^4,
    \qquad
    R_1=g_0^1g_1^3.
\]
Let the verifier challenge be
\[
    r_1=7,
    \qquad
    r_1^{-1}=14\pmod{97}.
\]
Fold the coefficients:
\[
    f'_0=r_1f_0+f_2=7\cdot 2+1=15,
\]
\[
    f'_1=r_1f_1+f_3=7\cdot 4+3=31.
\]
So
\[
    \mathbf f'=(15,31).
\]
Fold the bases:
\[
    g'_0=g_0^{14}g_2,
    \qquad
    g'_1=g_1^{14}g_3.
\]
Fold the value:
\[
    v'=r_1v_L+v_R=7\cdot22+16=170\equiv 73\pmod{97}.
\]
Indeed,
\[
    f'(5)=15+31\cdot5=170\equiv 73.
\]

Now check the commitment update:
\[
    \mathsf{com}'=L_1^{7}\mathsf{com}_fR_1^{14}.
\]
Expanding:
\[
\begin{aligned}
\mathsf{com}'
&=(g_2^2g_3^4)^7(g_0^2g_1^4g_2g_3^3)(g_0g_1^3)^{14}\\
&=g_0^{2+14}g_1^{4+42}g_2^{14+1}g_3^{28+3}\\
&=g_0^{16}g_1^{46}g_2^{15}g_3^{31}.
\end{aligned}
\]
On the other hand,
\[
\begin{aligned}
(g'_0)^{15}(g'_1)^{31}
&=(g_0^{14}g_2)^{15}(g_1^{14}g_3)^{31}\\
&=g_0^{210}g_2^{15}g_1^{434}g_3^{31}\\
&=g_0^{16}g_1^{46}g_2^{15}g_3^{31}\pmod{97},
\end{aligned}
\]
because $210\equiv16$ and $434\equiv46$ modulo $97$. Thus the folded commitment is correct.

\subsection{Second folding round and final scalar}

The remaining vector has length $2$. Split
\[
    \mathbf f'=(15,31)
\]
into left scalar $15$ and right scalar $31$. The verifier first checks the value decomposition
\[
    v'=15+u\cdot31=15+5\cdot31=170\equiv73.
\]
The prover sends
\[
    L_2=(g'_1)^{15},
    \qquad
    R_2=(g'_0)^{31}.
\]
Let the second challenge be
\[
    r_2=11,
    \qquad
    r_2^{-1}=53\pmod{97}.
\]
The final folded scalar is
\[
    f''=r_2\cdot15+31=196\equiv2\pmod{97}.
\]
The final folded base is
\[
    g''=(g'_0)^{53}g'_1.
\]
The second commitment update is
\[
    \mathsf{com}''=L_2^{r_2}\mathsf{com}'R_2^{r_2^{-1}}.
\]
Expanding with $\mathsf{com}'=(g'_0)^{15}(g'_1)^{31}$ gives
\[
\begin{aligned}
\mathsf{com}''
&=((g'_1)^{15})^{11}(g'_0)^{15}(g'_1)^{31}((g'_0)^{31})^{53}\
&=(g'_0)^{15+31\cdot 53}(g'_1)^{31+15\cdot 11}.
\end{aligned}
\]
Modulo $97$,
\[
    15+31\cdot53=1658\equiv 9,
    \qquad
    31+15\cdot11=196\equiv2.
\]
On the other hand,
\[
\begin{aligned}
(g'')^{f''}
&=((g'_0)^{53}g'_1)^2\
&=(g'_0)^{106}(g'_1)^2\
&=(g'_0)^9(g'_1)^2,
\end{aligned}
\]
since $106\equiv9\pmod{97}$. Thus
\[
    \mathsf{com}''=(g'')^{f''},
\]
so the commitment relation survives the second folding round as well.

This is the whole logic of Bulletproofs-style opening: a long vector commitment and inner-product claim are recursively folded into a scalar claim. The proof length grows only logarithmically in the original vector length.

